\section{Evaluation Setup}\label{sec:setup}

\subsection{Domain and Knowledge Source}

The evaluation domain is the Bank of Korea Compensation Regulations
(\emph{보수규정 전문, 2025-02-13 revision}), a publicly available
governance document of approximately 200 pages covering job grades,
positions, salary tables (별표 1--9), evaluation-tier differentials,
overseas pay, wage-peak schedules, and supplementary clauses (부칙).
The document mixes prose articles (제2조 definitions, 제5조 step
increments, 제9조 overtime, etc.) with multiple multi-key lookup tables.
It is small enough that the entire structured-fact set fits within a
32K-token context window, yet rich enough to expose multi-hop and
boundary reasoning.

\subsection{Question Set}

We construct a 100-question evaluation set in two stages. Thirty seed
questions were authored manually with verified gold answers traced to
the original document. Seventy additional questions were generated by
walking the underlying lookup tables with category-specific templates
(\texttt{eval/gen\_questions.py}), with weighted emphasis on multi-key
categories that stress the H1 hypothesis. The final distribution is
shown in Table~\ref{tab:eval-set}.

% Table: Evaluation set distribution (10 categories x 100 questions).
\begin{table}[t]
\centering
\caption{Evaluation set: 100 questions across ten categories. Multi-key
categories (calculation, comparison, multihop\_3key) total 42 questions
to ensure statistical power for the H1 comparison.}
\label{tab:eval-set}
\small
\begin{tabular}{l r r r}
\toprule
Category & Seed & Auto-expanded & Total \\
\midrule
single\_lookup        & 3 & 7 & 10 \\
join\_2key            & 3 & 7 & 10 \\
\textbf{multihop\_3key} & 3 & 12 & \textbf{15} \\
\textbf{calculation}    & 4 & 10 & \textbf{14} \\
\textbf{comparison}     & 3 & 10 & \textbf{13} \\
aggregation           & 3 & 7 & 10 \\
article\_text         & 3 & 5 & 8 \\
negation\_exception   & 3 & 5 & 8 \\
out\_of\_scope        & 2 & 4 & 6 \\
term\_normalization   & 3 & 3 & 6 \\
\midrule
\textbf{Total}        & \textbf{30} & \textbf{70} & \textbf{100} \\
\bottomrule
\end{tabular}
\end{table}


Each question carries a deterministic gold answer, an evidence pointer
to the source table or article, and a \texttt{required\_keys} list (a
proxy for join arity). Categories include single-key lookups
(\texttt{single\_lookup}), 2- and 3-key joins, calculations requiring
both lookup and arithmetic, comparisons across rows, aggregations,
exception clauses (negation), out-of-scope queries (combinations not
present in any table), and pure article-text questions.

\subsection{Backends}

We evaluate eight backends grouped into three classes:

\paragraph{Production agents.} The \backendname{typedb} and
\backendname{neo4j} backends invoke the production retrieval-augmented
agent: a LangGraph ReAct state machine in which a chat model performs
reasoning while a separate query model generates TypeQL or Cypher
queries against a live database. The \backendname{context} backend
uses production's chunked retrieval over a preprocessed regulation
document. We refer to this same backend as \backendname{C-MD} in the
surface-form Pareto comparison (Table~\ref{tab:token-cost} and
Figure~\ref{fig:pareto}), where it is the chunked-retrieval baseline
against the four full-dump in-context tiers. \backendname{closed\_book}
is the chat model with no retrieval — a baseline for measuring
pre-training leakage.

\paragraph{In-context surface-form ablation.} The remaining four
backends — \backendname{C-TQL}, \backendname{C-Cypher},
\backendname{C-NL}, and \backendname{C-JSONLD} — feed the \emph{entire}
knowledge graph to the chat model as a single in-context prompt, with
the only variable being the surface form. This decouples the
contribution of the query-generation stage from the in-context
reading capacity of the LLM. All four share the same article-text
prefix (the prose portion of the production context document) and
differ only in the structured-data section: TypeQL \texttt{insert}
statements (32{,}436 chars), Cypher \texttt{CREATE} statements
(32{,}030 chars), natural-language sentences (14{,}000 chars), or
JSON-LD-style nested objects (26{,}278 chars).

\subsection{Setting B: Decoupled Query Agent}

Pilot experiments revealed that when a single LLM serves both as
reasoner and as query generator, the KG agent's accuracy is dominated
by the LLM's TypeQL/Cypher syntax familiarity rather than by KG
expressivity. To remove this confound, all main experiments use
\emph{Setting B}: the chat model varies across the four LLM families
under test, while the query model is fixed to
\texttt{Qwen3-Coder-30B-A3B-Instruct}, a specialized code-generation
model that was the strongest available coder LLM in our gateway.
Production's \texttt{create\_chat\_model()} and
\texttt{create\_qwen\_model()} factories, originally designed for this
separation, are wired through environment variables.

\subsection{LLM Families}

Four LLM families serve as chat (reasoning) models. All are accessed
through a single OpenAI-compatible gateway with identical inference
parameters (temperature 0, no reasoning-mode flags); only the model
identifier varies.

\begin{itemize}
\item \textbf{HyperCLOVA X SEED-32B-Think-Text} — Korean-specialist
      think-mode model from NAVER. We confirm that without explicit
      reasoning-flag activation, no \texttt{<think>} traces appear in
      responses; HCX runs in non-reasoning mode for parity.
\item \textbf{Qwen3-235B-A22B-Instruct-2507-FP8} — frontier MoE
      instruct model from Alibaba (235B parameters, 22B active).
\item \textbf{Llama-4-Scout-17B-16E-Instruct} — Meta open-weight
      Western frontier (17B$\times$16 expert MoE).
\item \textbf{gpt-oss-120b} — OpenAI's 2025 open-weight release,
      serving as a stand-in for the closed GPT-4 family.
\end{itemize}

The query model (Qwen3-Coder-30B) is held constant across all KG-agent
runs.

\subsection{Metrics}

\paragraph{Per-item correctness.} Our scorer (\texttt{eval/scorer.py},
stdlib-only) handles seven gold-answer types: monetary amounts (with
currency and period), numbers, ratios (allowing percentage form),
booleans, sets, strings (with aliases), and refusals. For monetary
answers the scorer compares the closest numeric token in the
prediction against the gold value and verifies currency.

\paragraph{Aggregate metrics.}
For each (backend, LLM) pair we report point accuracy, the 95\%
percentile-bootstrap confidence interval (2{,}000 resamples), and a
paired McNemar exact two-sided binomial test on discordant pairs
against the relevant baseline (typically \backendname{context} or
\backendname{closed\_book}). All tests are paired by question id.

\paragraph{Surface-form size.} For the in-context surface forms, we
report the literal character count of the dump (the same string fed
to every prompt). Token cost scales roughly linearly with character
count for these largely numeric, semi-structured texts.
